% =====================================================================
%  CSE499A — EHR Pre-Consultation Medical Documentation System
%  Standalone section: Related Works
%  Authors: Rafiur Rahman Mashrafi (2221971042),
%           M.G. Rabbi Hossen (2222516042),
%           Israt Zaman Srity (2211084042)
%  Advisor: Dr. Mohammad Ashrafuzzaman Khan [AzK]
%  Department of Electrical and Computer Engineering,
%  North South University · Spring 2026
% =====================================================================

\documentclass[12pt,a4paper]{article}

\usepackage[utf8]{inputenc}
\usepackage[T1]{fontenc}
\usepackage{mathptmx}
\usepackage[a4paper,margin=1in]{geometry}
\usepackage{setspace}
\usepackage{enumitem}
\usepackage[hidelinks]{hyperref}
\usepackage{titlesec}
\usepackage{xcolor}

\onehalfspacing
\setlength{\parskip}{6pt}
\setlength{\parindent}{1.5em}

\titleformat{\section}{\normalfont\large\bfseries}{\thesection}{0.7em}{}
\titleformat{\subsection}{\normalfont\normalsize\bfseries}{\thesubsection}{0.7em}{}

\title{\bfseries Related Works:\\
EHR-Based Pre-Consultation Medical Documentation System}
\author{Rafiur Rahman Mashrafi (ID 2221971042)\\
        M.~G.~Rabbi Hossen (ID 2222516042)\\
        Israt Zaman Srity (ID 2211084042)\\[6pt]
        \emph{Faculty Advisor:} Dr. Mohammad Ashrafuzzaman Khan [AzK]\\
        Department of Electrical and Computer Engineering\\
        North South University, Bangladesh}
\date{Senior Design Project (CSE499A) \quad Spring 2026}

\begin{document}
\maketitle

\section{Scope of the Review}

The proposed system is a voice-driven pipeline that turns a
patient's spoken description --- in dialectal, code-mixed Bangla
--- into a structured Electronic Health Record (EHR) for the
attending physician. Because the pipeline integrates work from
several research streams, this related-works review is organized
into seven sub-areas: (i)~multilingual and large-scale automatic
speech recognition (ASR), (ii)~the recent generation of multimodal
audio language models, (iii)~Bangla-specific speech resources and
benchmarks, (iv)~clinical and biomedical natural language
processing, (v)~commercial AI clinical scribes and existing
voice-driven EHR products, (vi)~parameter-efficient fine-tuning
and audio-domain adaptation, and (vii)~AI ethics and governance for
clinical machine learning. Each subsection identifies how the
present project differs from, builds on, or extends the cited work.

\section{Multilingual and Large-Scale ASR}

Modern speech recognition is dominated by large self-supervised and
weakly supervised neural models trained on web-scale multilingual
audio. Radford et al.~\cite{rw:whisper} introduced
\textbf{Whisper}, an encoder--decoder Transformer trained on
$\sim$680{,}000 hours of weakly labeled multilingual audio that
achieves strong zero-shot performance across nearly 100 languages
including Bangla. Pratap et al.~\cite{rw:mms} released the
\textbf{Massively Multilingual Speech (MMS)} project, which extends
ASR coverage to over 1{,}100 languages by combining religious-text
audio with self-supervised pre-training and CTC fine-tuning. Baevski
et al.~\cite{rw:wav2vec2} proposed \textbf{wav2vec~2.0}, a
self-supervised speech-representation framework, and Conneau et
al.~\cite{rw:xlsr} extended it cross-lingually as XLSR-53;
both underpin many open Bangla ASR variants. The Seamless
Communication team's \textbf{SeamlessM4T}~\cite{rw:seamless}
unifies speech-to-text, text-to-text, and speech-to-speech
translation across a wide range of languages.

The original Transformer architecture~\cite{rw:vaswani} provides the
backbone of all of these models, while attention-mechanism
extensions for long-form audio (e.g.\ chunked attention) have
addressed the practical 30-second context limit that constrains
naive Whisper deployment. None of the above papers specifically
report results on dialectal Bangla clinical speech, which is the
precise gap the present project investigates.

\section{Multimodal Audio Language Models}

A recent direction fuses speech encoders with autoregressive
language-model decoders to create a single unified audio-LLM.
Chu et al.~\cite{rw:qwen2audio} introduced the \textbf{Qwen2-Audio}
family, and the same group followed with the specialized
\textbf{Qwen3-ASR-1.7B} variant. Mistral AI's
\textbf{Voxtral-Mini}~\cite{rw:voxtral} and Microsoft's
\textbf{Phi-4 Multimodal}~\cite{rw:phi4} extend this direction.
While these models report strong multilingual benchmarks on read
speech, their out-of-distribution behavior on dialect-rich
low-resource speech has been studied relatively little. The
empirical contribution of CSE499A is precisely on this point: a
controlled comparison showed that on real Bangladeshi dialect
audio, the smaller Bangla-specialized BengaliAI Regional model
outperformed all six 1.7B--7B parameter audio LLMs by more than
50 percentage points of WER --- a finding that does not appear in
any of the cited papers.

\section{Bangla Speech Resources and Benchmarks}

For Bengali specifically, several public resources exist.
The \textbf{Bengali.AI Speech Recognition
Competition}~\cite{rw:bengaliai} on Kaggle (2023) produced a
read-aloud Bengali benchmark and motivated several fine-tuned
Whisper and Wav2Vec~2.0 variants used in this project's baseline.
Mozilla's \textbf{Common Voice}~\cite{rw:commonvoice} provides a
crowdsourced Bengali corpus collected from volunteer speakers
worldwide. AI4Bharat's \textbf{IndicWav2Vec}~\cite{rw:indicwav2vec}
extends self-supervised speech models to twenty-two Indic
languages, and the \textbf{Vakyansh}~\cite{rw:vakyansh} project
releases ASR models tuned for Indic ASR including Bengali. The
\textbf{Shrutilipi} dataset~\cite{rw:shrutilipi} provides
$\sim$6{,}400 hours of paired audio--text data across twelve Indic
languages and is one of the largest open Bangla speech resources.

Most relevantly, Rakib et al.~\cite{rw:oodspeech} released the
\textbf{OOD-Speech} dataset, an out-of-distribution Bengali ASR
corpus that explicitly tests model robustness on speech that differs
from the training distribution. OOD-Speech is the closest published
work to the present project's dialect-aware focus. The two projects
differ in two important respects: OOD-Speech examines OOD-ness in
general (speaker, recording condition, register), while the current
project specifically targets the simultaneous combination of
\emph{dialect}, \emph{Bangla--English code-mixing}, and
\emph{clinical/informal register}. This intersection has not been
benchmarked in published work to date.

\section{Clinical and Biomedical NLP}

On the language-modeling side, Bhattacharjee et
al.~\cite{rw:banglabert} introduced \textbf{BanglaBERT}, a Bangla
BERT-style language model with benchmarks for multiple downstream
Bangla NLP tasks; this is the most likely starting point for the
medical NER work planned for CSE499B. For the medical domain, Lee
et al.~\cite{rw:biobert} released \textbf{BioBERT}, a BERT model
pretrained on PubMed abstracts and PMC full-text articles, and
Alsentzer et al.~\cite{rw:clinicalbert} released
\textbf{ClinicalBERT}, pretrained on MIMIC-III clinical notes;
both established that domain-adaptive pretraining substantially
improves downstream clinical extraction tasks.

Domain-specific clinical NLP also includes
\textbf{SciBERT}~\cite{rw:scibert} for scientific text and
\textbf{BlueBERT}~\cite{rw:bluebert}, which combines PubMed and
MIMIC-III pretraining. The original Transformer
architecture~\cite{rw:vaswani} and the BERT
encoder~\cite{rw:bert} underpin all of the above. The
Hugging~Face Transformers library~\cite{rw:transformers}
provides reference implementations and is the primary tooling used
in the present project.

A clear gap exists for Bangla: there is no published Bangla
equivalent of BioBERT or ClinicalBERT, no public Bangla medical NER
dataset, and no published clinical extraction benchmark in Bangla.
The CSE499B medical-NER component of this project, when delivered,
will be a small first step in this direction.

\section{Commercial AI Clinical Scribes and Voice EHR Systems}

The closest commercial analogues to the proposed system are the
recent generation of ``ambient AI scribes,'' which listen to a
doctor--patient conversation and generate a structured note. The
market-leading systems are \textbf{Nuance DAX
Copilot}~\cite{rw:nuance}, \textbf{DeepScribe}~\cite{rw:deepscribe},
\textbf{Abridge}~\cite{rw:abridge}, and \textbf{Suki AI}. Common
across these products are: (i)~English-only language coverage,
(ii)~deployment in well-resourced Western hospital systems, and
(iii)~tight integration with U.S.\ EHR platforms (Epic, Cerner).
None of them addresses Bangla, and none is designed for the
constraints of small Bangladeshi clinics: intermittent connectivity,
cost sensitivity, low patient digital literacy, and the absence of
a baseline EHR system to integrate against.

A second, older line of work concerns voice-driven EHR data entry
by physicians (e.g.\ Dragon Medical One). These systems treat the
doctor as the speaker, dictating after the visit. The present
project differs structurally: the \emph{patient} is the speaker,
during the wait, and the artifact produced is the \emph{intake
record} consulted by the physician at the start of the consultation.

\section{Parameter-Efficient Fine-Tuning and Domain Adaptation}

Hu et al.~\cite{rw:lora} introduced \textbf{Low-Rank
Adaptation (LoRA)}, which freezes the base model's weights and
trains only small low-rank adapter matrices, dramatically reducing
the compute and memory footprint of fine-tuning. LoRA and related
parameter-efficient fine-tuning (PEFT) techniques are now the
standard approach for adapting large speech and language models to
low-resource domains, and are the planned methodology for the
CSE499B fine-tuning phase of this project.

Houlsby et al.~\cite{rw:adapters} originally introduced adapter
modules in NLP and motivated the idea of inserting small trainable
layers into a frozen pretrained model. Bain et
al.'s~\cite{rw:whisperx} \textbf{WhisperX} provides robust
long-form audio transcription with forced alignment and
voice-activity detection; it is used in the audio-preprocessing
pipeline of the current project to segment long audio recordings
into the 10--15-second clips that the evaluated ASR models can
accept. Bredin et al.~\cite{rw:pyannote} released the
\textbf{pyannote.audio} library for speaker diarization and
voice-activity detection, which is the upstream of WhisperX's
alignment.

\section{AI Ethics and Governance for Clinical ML}

Because the system handles sensitive health data, prior work in
clinical AI ethics is directly relevant. The
\textbf{ACM Code of Ethics and Professional
Conduct}~\cite{rw:acm2018} provides the most concrete professional
guidance for software-systems work and has shaped the project's
patient-consent and transparency design. Beauchamp and Childress's
canonical \textbf{Principles of Biomedical
Ethics}~\cite{rw:beauchamp} (autonomy, non-maleficence, beneficence,
justice) is the framework physicians themselves use, and the
project's choice to keep the doctor as the final decision-maker is
grounded in patient autonomy and the principle of non-maleficence.
The \textbf{WHO guidance on the ethics and governance of AI for
health}~\cite{rw:whoethics} (2021) lays out six guiding
principles that have shaped the project's framing of human
oversight, transparency, and equitable access.

For a treatment of the broader risk surface introduced by clinical
ML, Char et al.~\cite{rw:char} and Vayena et
al.~\cite{rw:vayena} are widely cited surveys.

\section{Positioning of the Present Project}

The literature surveyed above shows that:

\begin{enumerate}[leftmargin=*]
\item There is mature open-source ASR for many high-resource
languages, but \emph{no benchmark for dialect-aware Bangla clinical
speech}. The Phase~2 results of this project (twelve baseline
models tested on a 4.7-hour multi-dialect corpus) help close that
gap.
\item Larger multimodal audio LLMs offer impressive generic
capability, but their behavior on dialect-rich low-resource
speech is largely unstudied. The Phase~3 results of this project
(six audio LLMs tested on the same corpus) offer a controlled
empirical observation that scaling alone does not solve dialect
Bangla ASR.
\item Clinical NLP infrastructure (BioBERT, ClinicalBERT,
ambient scribes) exists for English but \emph{does not exist for
Bangla}. The proposed end-to-end pipeline contributes a first step
toward Bangla clinical NLP infrastructure.
\item Parameter-efficient fine-tuning is the established way to
adapt large speech models to a new domain on a limited compute
budget, and is the planned approach for CSE499B.
\item Existing commercial ambient AI scribes target English-language
Western hospital workflows and do not transfer to small
Bangladeshi clinics, motivating a new design rather than a
re-skinning of existing products.
\end{enumerate}

To the best of the team's knowledge, no published end-to-end
voice-to-EHR system targets Bangla clinical use, and no published
benchmark addresses dialect, code-mixing, and informal clinical
register simultaneously. These two observations define the
contribution that the project aims to make over CSE499A and
CSE499B.

% =====================================================================
% REFERENCES
% =====================================================================
\begin{thebibliography}{99}

\bibitem{rw:whisper}
A.~Radford, J.~W.~Kim, T.~Xu, G.~Brockman, C.~McLeavey,
and I.~Sutskever,
``Robust Speech Recognition via Large-Scale Weak Supervision,''
in \emph{Proc.\ 40th Int.\ Conf.\ on Machine Learning (ICML)},
pp.~28492--28518, 2023.
\url{https://arxiv.org/abs/2212.04356}.

\bibitem{rw:mms}
V.~Pratap \emph{et al.},
``Scaling Speech Technology to 1{,}000+ Languages,''
\emph{arXiv preprint arXiv:2305.13516}, 2023.
\url{https://arxiv.org/abs/2305.13516}.

\bibitem{rw:wav2vec2}
A.~Baevski, Y.~Zhou, A.~Mohamed, and M.~Auli,
``wav2vec 2.0: A Framework for Self-Supervised Learning of Speech
Representations,''
in \emph{Advances in Neural Information Processing Systems
(NeurIPS)}, vol.~33, pp.~12449--12460, 2020.
\url{https://arxiv.org/abs/2006.11477}.

\bibitem{rw:xlsr}
A.~Conneau, A.~Baevski, R.~Collobert, A.~Mohamed, and M.~Auli,
``Unsupervised Cross-Lingual Representation Learning for Speech
Recognition,''
in \emph{Proc.\ INTERSPEECH 2021}, pp.~2426--2430.
\url{https://arxiv.org/abs/2006.13979}.

\bibitem{rw:seamless}
Seamless Communication Team \emph{et al.},
``SeamlessM4T: Massively Multilingual and Multimodal Machine
Translation,''
\emph{arXiv preprint arXiv:2308.11596}, 2023.
\url{https://arxiv.org/abs/2308.11596}.

\bibitem{rw:vaswani}
A.~Vaswani \emph{et al.},
``Attention Is All You Need,''
in \emph{Advances in Neural Information Processing Systems
(NeurIPS)}, vol.~30, 2017.
\url{https://arxiv.org/abs/1706.03762}.

\bibitem{rw:qwen2audio}
Y.~Chu \emph{et al.},
``Qwen2-Audio Technical Report,''
\emph{arXiv preprint arXiv:2407.10759}, 2024.
\url{https://arxiv.org/abs/2407.10759}.

\bibitem{rw:voxtral}
Mistral AI,
``Voxtral: Open Multilingual Speech-to-Text Models.''
\url{https://mistral.ai/news/voxtral/},
2024. Online; accessed April 2026.

\bibitem{rw:phi4}
Microsoft Research,
``Phi-4 Multimodal: Compact Multimodal Foundation Models,''
2024. \url{https://huggingface.co/microsoft/Phi-4-multimodal-instruct}.

\bibitem{rw:bengaliai}
Bengali.AI,
``Bengali.AI Speech Recognition Competition,''
Kaggle, 2023.
\url{https://www.kaggle.com/competitions/bengaliai-speech}.

\bibitem{rw:commonvoice}
R.~Ardila \emph{et al.},
``Common Voice: A Massively-Multilingual Speech Corpus,''
in \emph{Proc.\ 12th Language Resources and Evaluation Conf.\
(LREC)}, pp.~4218--4222, 2020.
\url{https://commonvoice.mozilla.org}.

\bibitem{rw:indicwav2vec}
AI4Bharat,
``IndicWav2Vec: Self-Supervised Speech Models for Indian Languages,''
2022. \url{https://github.com/AI4Bharat/IndicWav2Vec}.

\bibitem{rw:vakyansh}
A.~Gupta \emph{et al.},
``Vakyansh: ASR Toolkit for Low Resource Indic Languages,''
2022. \url{https://github.com/Open-Speech-EkStep/vakyansh-models}.

\bibitem{rw:shrutilipi}
K.~Bhogale \emph{et al.},
``Effectiveness of Mining Audio and Text Pairs from Public Data for
Improving ASR Systems for Low-Resource Languages,''
in \emph{Proc.\ ICASSP 2023}.
\url{https://arxiv.org/abs/2208.12666}.

\bibitem{rw:oodspeech}
F.~R.~Rakib \emph{et al.},
``OOD-Speech: A Large Bengali Speech Recognition Dataset for
Out-of-Distribution Benchmarking,''
in \emph{Proc.\ INTERSPEECH 2023}.
\url{https://arxiv.org/abs/2305.09688}.

\bibitem{rw:banglabert}
A.~Bhattacharjee \emph{et al.},
``BanglaBERT: Language Model Pretraining and Benchmarks for
Low-Resource Language Understanding Evaluation in Bangla,''
in \emph{Findings of NAACL 2022}, pp.~1318--1327.
\url{https://github.com/csebuetnlp/banglabert}.

\bibitem{rw:biobert}
J.~Lee \emph{et al.},
``BioBERT: A Pre-Trained Biomedical Language Representation Model
for Biomedical Text Mining,''
\emph{Bioinformatics}, vol.~36, no.~4, pp.~1234--1240, 2020.
\url{https://arxiv.org/abs/1901.08746}.

\bibitem{rw:clinicalbert}
E.~Alsentzer \emph{et al.},
``Publicly Available Clinical BERT Embeddings,''
in \emph{Proc.\ 2nd Clinical NLP Workshop, NAACL 2019},
pp.~72--78.
\url{https://arxiv.org/abs/1904.03323}.

\bibitem{rw:scibert}
I.~Beltagy, K.~Lo, and A.~Cohan,
``SciBERT: A Pretrained Language Model for Scientific Text,''
in \emph{Proc.\ EMNLP-IJCNLP 2019}.
\url{https://arxiv.org/abs/1903.10676}.

\bibitem{rw:bluebert}
Y.~Peng, S.~Yan, and Z.~Lu,
``Transfer Learning in Biomedical Natural Language Processing:
An Evaluation of BERT and ELMo on Ten Benchmarking Datasets,''
in \emph{Proc.\ BioNLP 2019 Workshop}.
\url{https://arxiv.org/abs/1906.05474}.

\bibitem{rw:bert}
J.~Devlin, M.-W.~Chang, K.~Lee, and K.~Toutanova,
``BERT: Pre-training of Deep Bidirectional Transformers for
Language Understanding,''
in \emph{Proc.\ NAACL-HLT 2019}, pp.~4171--4186.
\url{https://arxiv.org/abs/1810.04805}.

\bibitem{rw:transformers}
T.~Wolf \emph{et al.},
``Transformers: State-of-the-Art Natural Language Processing,''
in \emph{Proc.\ EMNLP 2020: System Demonstrations}, pp.~38--45.
\url{https://github.com/huggingface/transformers}.

\bibitem{rw:nuance}
Microsoft / Nuance Communications,
``Nuance DAX Copilot --- Ambient AI Documentation.''
\url{https://www.nuance.com/healthcare/dragon-ai-clinical-solutions/dax-copilot.html}.
Online; accessed April 2026.

\bibitem{rw:deepscribe}
DeepScribe Inc.,
``DeepScribe AI Medical Scribe.''
\url{https://www.deepscribe.ai}. Online; accessed April 2026.

\bibitem{rw:abridge}
Abridge AI Inc.,
``Abridge: Generative AI for Clinical Conversations.''
\url{https://www.abridge.com}. Online; accessed April 2026.

\bibitem{rw:lora}
E.~J.~Hu \emph{et al.},
``LoRA: Low-Rank Adaptation of Large Language Models,''
in \emph{Int.\ Conf.\ on Learning Representations (ICLR)}, 2022.
\url{https://arxiv.org/abs/2106.09685}.

\bibitem{rw:adapters}
N.~Houlsby \emph{et al.},
``Parameter-Efficient Transfer Learning for NLP,''
in \emph{Proc.\ ICML 2019}.
\url{https://arxiv.org/abs/1902.00751}.

\bibitem{rw:whisperx}
M.~Bain, J.~Huh, T.~Han, and A.~Zisserman,
``WhisperX: Time-Accurate Speech Transcription of Long-Form Audio,''
in \emph{Proc.\ INTERSPEECH 2023}.
\url{https://arxiv.org/abs/2303.00747}.

\bibitem{rw:pyannote}
H.~Bredin \emph{et al.},
``pyannote.audio: Neural Building Blocks for Speaker Diarization,''
in \emph{Proc.\ ICASSP 2020}.
\url{https://github.com/pyannote/pyannote-audio}.

\bibitem{rw:acm2018}
ACM Council,
``ACM Code of Ethics and Professional Conduct,'' 2018.
\url{https://www.acm.org/code-of-ethics}.

\bibitem{rw:beauchamp}
T.~L.~Beauchamp and J.~F.~Childress,
\emph{Principles of Biomedical Ethics}, 8th ed.,
Oxford University Press, 2019.

\bibitem{rw:whoethics}
World Health Organization,
``Ethics and Governance of Artificial Intelligence for Health,''
WHO Guidance, 2021.
\url{https://www.who.int/publications/i/item/9789240029200}.

\bibitem{rw:char}
D.~S.~Char, N.~H.~Shah, and D.~Magnus,
``Implementing Machine Learning in Health Care --- Addressing
Ethical Challenges,''
\emph{New England Journal of Medicine}, vol.~378, pp.~981--983,
2018.

\bibitem{rw:vayena}
E.~Vayena, A.~Blasimme, and I.~G.~Cohen,
``Machine Learning in Medicine: Addressing Ethical Challenges,''
\emph{PLOS Medicine}, vol.~15, no.~11, e1002689, 2018.

\end{thebibliography}

\end{document}
